\documentclass[11pt,a4paper]{article}
\usepackage[top=1.6cm,bottom=1.4cm,left=1.3cm,right=1.3cm,headsep=8pt,footskip=16pt]{geometry}
\usepackage{enumitem}
\usepackage{amssymb}
\usepackage{array}
\usepackage{xcolor}
\usepackage{fancyhdr}
\usepackage{multicol}
\usepackage{polyglossia}
\setdefaultlanguage[calendar=gregorian, numerals=maghrib]{arabic}
\setotherlanguage{english}
\IfFontExistsTF{Amiri}{%
  \newfontfamily\arabicfont[Script=Arabic, Scale=1.08]{Amiri}
}{%
  \newfontfamily\arabicfont[Script=Arabic, Scale=1.08, Path=../fonts/, Extension=.ttf,
    UprightFont=*-Regular, BoldFont=*-Bold, ItalicFont=*-Italic,
    BoldItalicFont=*-BoldItalic]{Amiri}
}
\IfFontExistsTF{AlBattar}{%
  \newfontfamily\khatt[Script=Arabic, Scale=1.15]{AlBattar}
}{%
  \newfontfamily\khatt[Script=Arabic, Scale=1.15, Path=../fonts/,
    BoldFont=ae_AlBattar.ttf]{ae_AlBattar.ttf}
}
\renewcommand{\headrulewidth}{0.8pt}
\renewcommand{\footrulewidth}{0.8pt}
\fancyhf{}
\rhead{\khatt المدرسة العليا للأساتذة ببوزريعة}
\lhead{\khatt ملحقة بومرداس}
\lfoot{\khatt الأستاذ: د. بن عبد الرزاق يوسف}
\cfoot{}
\pagestyle{fancy}
\parindent 0pt
\setlength{\columnsep}{0.6cm}
\setlength{\columnseprule}{0.4pt}
\setlist[enumerate,1]{label=\textbf{\arabic*.}, leftmargin=*, itemsep=5pt, topsep=2pt}
\begin{document}
\begin{center}
{\Large\khatt\textbf{مفتاح الإجابات ــ الاختبارات من 1 إلى 12}}\\[4pt]
{\small (نسخة الأستاذ ــ لا تُوزَّع على الطلبة)}\\[10pt]
\renewcommand{\arraystretch}{1.8}
\begin{tabular}{|>{\centering\arraybackslash}m{2.2cm}|*{15}{>{\centering\arraybackslash}m{0.62cm}|}}
\hline
\textbf{رقم السؤال} & \textbf{1} & \textbf{2} & \textbf{3} & \textbf{4} & \textbf{5} & \textbf{6} & \textbf{7} & \textbf{8} & \textbf{9} & \textbf{10} & \textbf{11} & \textbf{12} & \textbf{13} & \textbf{14} & \textbf{15} \\ \hline\hline
\textbf{الاختبار 1} & ج & أ & ج & أ & د & هـ & د & د & هـ & هـ & ب & ج & أ & ب & ب \\ \hline
\textbf{الاختبار 2} & أ & هـ & ب & أ & د & ج & د & هـ & ب & ب & أ & هـ & ج & د & ج \\ \hline
\textbf{الاختبار 3} & د & ج & هـ & ب & د & هـ & هـ & ب & د & ب & ج & أ & ج & أ & أ \\ \hline
\textbf{الاختبار 4} & ب & ج & أ & ب & هـ & هـ & ج & ب & د & ج & د & أ & د & أ & هـ \\ \hline
\textbf{الاختبار 5} & أ & هـ & أ & ج & ج & ب & ج & هـ & د & د & أ & د & ب & ب & هـ \\ \hline
\textbf{الاختبار 6} & د & أ & د & ج & ب & هـ & ج & هـ & ج & هـ & أ & أ & ب & د & ب \\ \hline
\textbf{الاختبار 7} & أ & هـ & ج & هـ & د & د & ب & ب & ب & ج & أ & ج & أ & هـ & د \\ \hline
\textbf{الاختبار 8} & د & هـ & ب & ب & د & أ & هـ & د & ج & ب & ج & أ & هـ & أ & ج \\ \hline
\textbf{الاختبار 9} & ب & ب & أ & ج & هـ & د & أ & ب & د & د & هـ & هـ & أ & ج & ج \\ \hline
\textbf{الاختبار 10} & د & ب & أ & هـ & أ & ب & ب & ج & ج & أ & د & هـ & هـ & ج & د \\ \hline
\textbf{الاختبار 11} & ج & هـ & د & أ & ب & أ & د & ب & ج & هـ & د & ج & ب & أ & هـ \\ \hline
\textbf{الاختبار 12} & ب & د & أ & ج & أ & أ & هـ & د & ج & ب & هـ & هـ & ج & د & ب \\ \hline
\end{tabular}

\vspace{10pt}
{\small سلّم التنقيط: نقطة واحدة لكل سؤال (المجموع 15 نقطة). ترتيب الأسئلة: 1--3 سهلة، 4--10 متوسطة، 11--15 صعبة.}
\end{center}
\end{document}
